\documentclass[12pt,a4paper]{article}

\usepackage[utf8]{inputenc}
\usepackage[T1]{fontenc}
\usepackage[spanish]{babel}
\usepackage{amsmath, amssymb}
\usepackage{geometry}
\usepackage{booktabs}
\usepackage{array}
\usepackage{xcolor}
\usepackage{hyperref}
\usepackage{enumitem}
\usepackage{fancyhdr}
\usepackage{graphicx}
\usepackage{listings}
\usepackage{tcolorbox}
\usepackage{multirow}
\usepackage{caption}

\geometry{top=2.5cm, bottom=2.5cm, left=3cm, right=3cm}

% ── Colores ──────────────────────────────────────────────────────────────────
\definecolor{uddblue}{RGB}{0, 75, 135}
\definecolor{lightgray}{RGB}{245, 245, 245}
\definecolor{codegray}{RGB}{100, 100, 100}
\definecolor{codegreen}{RGB}{0, 120, 0}
\definecolor{codeblue}{RGB}{0, 0, 180}
\definecolor{codered}{RGB}{160, 0, 0}

% ── Estilo de código Python ───────────────────────────────────────────────────
\lstset{
    language     = Python,
    backgroundcolor = \color{lightgray},
    basicstyle   = \small\ttfamily,
    keywordstyle = \color{codeblue}\bfseries,
    commentstyle = \color{codegreen}\itshape,
    stringstyle  = \color{codered},
    numberstyle  = \tiny\color{codegray},
    numbers      = left,
    stepnumber   = 1,
    breaklines   = true,
    frame        = single,
    rulecolor    = \color{gray!40},
    tabsize      = 4,
    showstringspaces = false,
    literate =
        {á}{{\'{a}}}1 {é}{{\'{e}}}1 {í}{{\'{i}}}1
        {ó}{{\'{o}}}1 {ú}{{\'{u}}}1 {ñ}{{\~{n}}}1
        {Á}{{\'{A}}}1 {É}{{\'{E}}}1 {Í}{{\'{I}}}1
        {Ó}{{\'{O}}}1 {Ú}{{\'{U}}}1 {Ñ}{{\~{N}}}1
}

% ── Cajas de ayuda ────────────────────────────────────────────────────────────
\tcbuselibrary{skins, breakable}
\newtcolorbox{ayuda}[1]{
    colback   = lightgray,
    colframe  = uddblue,
    fonttitle = \bfseries\small,
    title     = {#1},
    breakable,
    left      = 6pt,
    right     = 6pt,
    top       = 4pt,
    bottom    = 4pt
}

% ── Encabezado y pie ─────────────────────────────────────────────────────────
\pagestyle{fancy}
\fancyhf{}
\fancyhead[L]{\small\color{gray} IIP314W\,--\,Optimización Aplicada a Negocios}
\fancyhead[R]{\small\color{gray} Tarea 4\,--\,2026-T2}
\fancyfoot[C]{\small\thepage}
\renewcommand{\headrulewidth}{0.4pt}

\hypersetup{colorlinks=true, linkcolor=uddblue, urlcolor=uddblue}

% ── Espaciado entre párrafos ──────────────────────────────────────────────────
\setlength{\parskip}{0.5em}
\setlength{\parindent}{0pt}

% =============================================================================
\begin{document}
% =============================================================================

% ── Título ───────────────────────────────────────────────────────────────────
\begin{center}
    {\Large\bfseries\color{uddblue}
     (IIP314W) Optimización Aplicada a Negocios}\\[0.4em]
    {\large Universidad del Desarrollo}\\[0.8em]
    \noindent\rule{\linewidth}{1.5pt}\\[0.5em]
    {\LARGE\bfseries Tarea 4}\\[0.3em]
    {\large Red de Agroexportación con Inventario Perecible:\\
     Modelamiento Multiperiodo, Dualidad y Holgura Complementaria}\\[0.5em]
    \noindent\rule{\linewidth}{1.5pt}\\[0.8em]
    \begin{tabular}{@{}ll@{}}
        \textbf{Profesor:} & Ing.\ Rodrigo Trigo Vilches \\[0.3em]
        \textbf{Ayudante:} & Lic.\ Vicente Ramírez Almonacid \\[0.3em]
        \textbf{Período:}  & Año 2026\,--\,Trimestre 2
    \end{tabular}
\end{center}

\vspace{1em}
\noindent\rule{\linewidth}{0.4pt}
\vspace{0.5em}

% =============================================================================
\section*{Contexto: AgroAndes Exporta S.A.}
% =============================================================================

\textbf{AgroAndes Exporta S.A.}\ es una empresa exportadora de fruta fresca que
debe planificar su red logística a lo largo de la \textbf{temporada de cosecha},
dividida en un conjunto de \textbf{periodos} (semanas) $t\in\mathcal{T}$. La red
tiene tres eslabones y la fruta fluye siempre en el mismo sentido:
\[
\text{\textbf{packing}} \;\longrightarrow\;
\text{\textbf{frigoríficos} (cámaras de frío)} \;\longrightarrow\;
\text{\textbf{mercados} de exportación}.
\]

En el primer eslabón están las plantas de \textbf{packing} $p\in\mathcal{P}$, que
reciben la fruta recién cosechada, la seleccionan y la embalan en pallets listos
para exportar. Aquí aparece la primera particularidad del negocio: \textbf{la
cosecha es estacional}. Solo durante las primeras semanas de la temporada
ingresa fruta fresca a proceso; después la temporada termina y ya \emph{no hay
más fruta que procesar}. Esto se refleja en que la capacidad de proceso
$\mathit{Cap}_{pt}$ de cada packing es positiva solo en los periodos de cosecha y
\textbf{cae a cero} una vez que la temporada cierra. Además, procesar una tonelada
de fruta en el packing $p$ durante el periodo $t$ tiene un \textbf{costo unitario}
$\mathit{cf}_{pt}$ (compra de fruta al productor más proceso) que \emph{cambia a lo
largo de la temporada}: al comienzo, con la cosecha en su punto máximo, la fruta
es abundante y barata; a medida que avanza la temporada se encarece.

En el segundo eslabón están los \textbf{frigoríficos} $b\in\mathcal{B}$, cámaras
de frío regionales que reciben la fruta embalada, la \textbf{almacenan} y desde
ahí despachan a los mercados. Como toda instalación, un frigorífico se puede
tener \textbf{operativo} o cerrado en cada periodo, y se puede \textbf{abrir} uno
que estaba cerrado o \textbf{cerrar} uno que estaba abierto, pagando cada vez el
costo de la transición. Aquí aparece la segunda —y central— particularidad: la
fruta almacenada es \textbf{perecible}. Una tonelada de fruta no puede permanecer
en cámara indefinidamente: pierde calidad de exportación y, pasado cierto tiempo,
ya no es despachable. Concretamente, \textbf{la fruta puede permanecer almacenada
a lo más $L$ periodos} desde que ingresa al frigorífico.

En el tercer eslabón están los \textbf{mercados} de exportación $j\in\mathcal{J}$,
cada uno con una \textbf{demanda} $d_{jt}$ (pedidos comprometidos) que debe
cumplirse en cada periodo.

El desajuste entre \emph{cuándo hay fruta} (temprano, barata, pero acotada por la
capacidad de cosecha) y \emph{cuándo hay demanda} (repartida a lo largo de toda la
temporada) obliga a la empresa a \textbf{almacenar}: producir fruta barata al
inicio y guardarla en frío para despacharla más adelante. Pero la perecibilidad
pone un límite duro a esa estrategia. La gerencia quiere el plan de
\emph{apertura, cierre, proceso, transporte y almacenamiento} que satisfaga toda
la demanda de la temporada al \textbf{mínimo costo total}.

\begin{ayuda}{Trabajo en grupo}
\vspace{2pt}
La tarea se desarrolla \textbf{en parejas o grupos de 3} (ni más ni menos). Tanto
el informe como el código deben incluir los nombres de todos los integrantes.
\textbf{Es obligatorio resolver el modelo con \texttt{gurobipy}}.
\end{ayuda}

\noindent\rule{\linewidth}{0.4pt}

% =============================================================================
\section*{Datos del problema}
% =============================================================================

La empresa opera con $|\mathcal{P}|=2$ plantas de packing ($A$, $B$),
$|\mathcal{B}|=3$ frigoríficos ($F_1, F_2, F_3$) y $|\mathcal{J}|=3$ mercados de
exportación ($X$, $Y$, $Z$), sobre un horizonte de $|\mathcal{T}|=4$ periodos
($t=1,2,3,4$). La \textbf{vida útil} de la fruta en cámara es $L=2$ periodos.
Todas las cantidades de fruta están en \textbf{toneladas}; todos los costos, en
\textbf{miles de USD} (costos unitarios en miles de USD por tonelada).

\begin{center}
\renewcommand{\arraystretch}{1.2}
\begin{tabular}{lcccc}
\toprule
\textbf{Demanda $d_{jt}$ (ton)} & \textbf{$t=1$} & \textbf{$t=2$} &
\textbf{$t=3$} & \textbf{$t=4$} \\
\midrule
Mercado $X$ & 25 & 35 & 45 & 55 \\
Mercado $Y$ & 15 & 20 & 25 & 30 \\
Mercado $Z$ & 0  & 5  & 10 & 15 \\
\midrule
\textbf{Total} & 40 & 60 & 80 & 100 \\
\bottomrule
\end{tabular}
\captionof{table}{Demanda comprometida por mercado y periodo.}
\end{center}

\textbf{Plantas de packing.} La capacidad de proceso $\mathit{Cap}_{pt}$ es
positiva solo durante la ventana de cosecha ($t=1,2$) y cae a cero después. El
costo unitario de proceso $\mathit{cf}_{pt}$ sube a lo largo de la temporada.

\begin{center}
\renewcommand{\arraystretch}{1.2}
\setlength{\tabcolsep}{5pt}
\begin{tabular}{lcccc@{\hskip 1.2em}cccc}
\toprule
 & \multicolumn{4}{c}{\textbf{$\mathit{Cap}_{pt}$ (ton)}} &
   \multicolumn{4}{c}{\textbf{$\mathit{cf}_{pt}$ (\$/ton)}} \\
\cmidrule(lr){2-5}\cmidrule(lr){6-9}
\textbf{Packing} & $t{=}1$ & $t{=}2$ & $t{=}3$ & $t{=}4$
                 & $t{=}1$ & $t{=}2$ & $t{=}3$ & $t{=}4$ \\
\midrule
$A$ & 100 & 100 & 0 & 0 & 4 & 9  & --- & --- \\
$B$ & 60  & 60  & 0 & 0 & 5 & 10 & --- & --- \\
\bottomrule
\end{tabular}
\captionof{table}{Capacidad de proceso y costo unitario de proceso por packing y
periodo. En $t=3,4$ no hay cosecha ($\mathit{Cap}=0$), por lo que $\mathit{cf}$ es
irrelevante.}
\end{center}

\begin{center}
\renewcommand{\arraystretch}{1.2}
\begin{tabular}{lcccc}
\toprule
\textbf{Packing} &
\textbf{$\mathit{FP}_{p}$ (fijo/per.)} &
\textbf{$\mathit{AP}_{p}$ (abrir)} &
\textbf{$\mathit{CP}_{p}$ (cerrar)} &
\textbf{Estado inicial $y_{p0}$} \\
\midrule
$A$ & 100 & 200 & 150 & Operativa (1) \\
$B$ & 90  & 150 & 120 & Cerrada (0) \\
\bottomrule
\end{tabular}
\captionof{table}{Costos fijos y de transición de las plantas de packing. El costo
fijo de operar se paga en cada periodo en que la planta está operativa.}
\end{center}

\textbf{Frigoríficos.} Cada frigorífico tiene una capacidad de almacenamiento
$\mathit{Alm}_b$, un costo fijo de operación $\mathit{FB}_b$ por periodo, costos de
apertura/cierre $\mathit{AB}_b,\mathit{CB}_b$ y un costo de mantención de
inventario $h_b$ por tonelada almacenada al cierre de cada periodo.

\begin{center}
\renewcommand{\arraystretch}{1.2}
\setlength{\tabcolsep}{6pt}
\begin{tabular}{lcccccc}
\toprule
\textbf{Frig.} & $\mathit{Alm}_{b}$ & $\mathit{FB}_{b}$ & $\mathit{AB}_{b}$ &
$\mathit{CB}_{b}$ & $h_{b}$ & \textbf{Estado inicial $z_{b0}$} \\
 & (ton) & (fijo/per.) & (abrir) & (cerrar) & (\$/ton) & \\
\midrule
$F_1$ & 80 & 40 & 120 & 80 & 1 & Operativo (1) \\
$F_2$ & 60 & 35 & 100 & 70 & 1 & Cerrado (0) \\
$F_3$ & 50 & 30 & 90  & 60 & 2 & Cerrado (0) \\
\bottomrule
\end{tabular}
\captionof{table}{Parámetros de los frigoríficos. El inventario inicial es
$s_{b0}=0$ para los tres.}
\end{center}

\textbf{Costos unitarios de transporte.} Enviar una tonelada de un packing $p$ a un
frigorífico $b$ cuesta $\mathit{cp}_{pb}$; enviarla de un frigorífico $b$ a un
mercado $j$ cuesta $\mathit{cg}_{bj}$.

\begin{center}
\renewcommand{\arraystretch}{1.2}
\begin{minipage}{0.46\linewidth}\centering
\begin{tabular}{lccc}
\toprule
\textbf{$\mathit{cp}_{pb}$} & $F_1$ & $F_2$ & $F_3$ \\
\midrule
$A$ & 3 & 4 & 6 \\
$B$ & 5 & 3 & 4 \\
\bottomrule
\end{tabular}
\captionof{table}{Packing $\to$ frigorífico.}
\end{minipage}\hfill
\begin{minipage}{0.46\linewidth}\centering
\begin{tabular}{lccc}
\toprule
\textbf{$\mathit{cg}_{bj}$} & $X$ & $Y$ & $Z$ \\
\midrule
$F_1$ & 2 & 5 & 8 \\
$F_2$ & 6 & 3 & 5 \\
$F_3$ & 8 & 6 & 3 \\
\bottomrule
\end{tabular}
\captionof{table}{Frigorífico $\to$ mercado.}
\end{minipage}
\end{center}

\noindent\textit{Nota sobre el big-M.} Para ligar el flujo entrante a un
frigorífico con su estado de apertura necesitarán una cota superior válida
(big-M) del flujo que puede entrar a un frigorífico en un periodo. Una cota
natural es la suma de las capacidades de proceso de los packing en ese periodo,
$M^{B}_{t}=\sum_{p}\mathit{Cap}_{pt}$. Recuerden elegir el big-M lo más ajustado
posible.

\noindent\rule{\linewidth}{0.4pt}

% =============================================================================
\section*{Parte A \quad Formulación del Modelo \hfill\normalfont\normalsize(30 pts)}
% =============================================================================

Formule un \textbf{modelo de optimización} que minimice el costo total de la
temporada. Su formulación debe ser completamente algebraica (en notación
matemática, \textbf{no} en código) y \textbf{paramétrica} (use la notación de los
conjuntos y parámetros, no los números concretos). \textbf{Toda la información
necesaria está en el Contexto y en los Datos}; una parte esencial de la tarea es
que \textbf{ustedes identifiquen} los componentes del modelo. Su formulación debe
contener:

\begin{enumerate}[label=\alph*), leftmargin=*, topsep=2pt]
    \item \textbf{Conjuntos e índices:} defina los conjuntos para indexar packing,
          frigoríficos, mercados y periodos.
    \item \textbf{Parámetros:} liste, con su notación, todos los datos del
          problema (incluyendo la vida útil $L$ y los estados iniciales).
    \item \textbf{Variables de decisión:} identifique y defina \textbf{todas} las
          variables. Necesitará binarias para el \textbf{estado} (operativo/cerrado)
          y para las \textbf{transiciones} (abrir/cerrar) de cada instalación, y
          variables continuas para los \textbf{flujos} en cada tramo y el
          \textbf{inventario}. Para cada una indique su \textbf{significado},
          \textbf{unidades}, \textbf{dominio} y \textbf{tipo}.
    \item \textbf{Función objetivo:} plantee el costo total a minimizar. Debe
          incluir: costos fijos de operación, costos de apertura y cierre, costo
          de proceso de la fruta ($\mathit{cf}_{pt}$), costos de transporte de
          ambos tramos y costo de inventario.
    \item \textbf{Restricciones:} formule y \textbf{nombre} todas las
          restricciones. Debe modelar, como mínimo:
          \begin{itemize}[leftmargin=1.4em, topsep=1pt]
            \item capacidad de proceso de cada packing, ligada a su estado
                  (un packing cerrado no procesa);
            \item activación del flujo entrante a un frigorífico según su estado
                  (un frigorífico cerrado no recibe fruta) y capacidad de
                  almacenamiento;
            \item \textbf{balance de inventario} en cada frigorífico y periodo
                  (lo que queda es lo que había, más lo que entró, menos lo que
                  salió);
            \item satisfacción de la demanda de cada mercado en cada periodo;
            \item lógica de \textbf{transición} abrir/cerrar que ligue el estado
                  entre periodos consecutivos (y con el estado inicial), impidiendo
                  abrir y cerrar simultáneamente;
            \item \textbf{perecibilidad}: la fruta no puede permanecer almacenada
                  más de $L$ periodos. \emph{Sugerencia:} el inventario al cierre
                  de un periodo $t$ solo puede estar compuesto por fruta que
                  ingresó al frigorífico en los últimos $L$ periodos.
          \end{itemize}
          Incluya los dominios de las variables.
\end{enumerate}

\noindent\textit{Sugerencia de método (no de solución): traduzca cada frase del
Contexto en una o más restricciones. La perecibilidad y la ventana de cosecha son
las dos ideas nuevas de esta tarea; el resto (costos fijos, aperturas/cierres,
balance de inventario) sigue la misma lógica de las ayudantías de modelamiento.}

\noindent\rule{\linewidth}{0.4pt}

% =============================================================================
\section*{Parte B \quad Implementación y Resolución en Gurobi \hfill\normalfont\normalsize(25 pts)}
% =============================================================================

Implemente el modelo de la Parte~A en \texttt{gurobipy} y resuélvalo con los datos
entregados. Responda:

\begin{enumerate}[label=\alph*), leftmargin=*, topsep=2pt]
    \item Reporte el \textbf{costo total óptimo} y descompóngalo por componente
          (fijos, transiciones, proceso, transporte tramo 1, transporte tramo 2,
          inventario).
    \item Presente el \textbf{plan de apertura y cierre}: en qué periodos está
          operativa cada planta y cada frigorífico, y en qué periodo exacto ocurre
          cada apertura o cierre.
    \item Presente la \textbf{trayectoria de inventario} $s_{bt}$ de cada
          frigorífico a lo largo del horizonte.
    \item Verifique e informe que se cumplen: (i) la cobertura de toda la demanda,
          (ii) el balance de inventario, y (iii) la restricción de perecibilidad
          (ninguna tonelada permanece almacenada más de $L$ periodos).
    \item Interprete la solución en \textbf{lenguaje de negocios}: ¿por qué se
          almacena fruta en los primeros periodos? ¿por qué se cierran las plantas
          de packing a mitad de temporada? ¿algún frigorífico queda sin usar?
\end{enumerate}

\begin{ayuda}{Ayuda -- Sintaxis de \texttt{gurobipy} (referencia del lenguaje)}
\vspace{2pt}
Esta caja es solo una referencia de \emph{sintaxis}, con nombres genéricos:
\textbf{no} es el modelo ni indica qué variables o restricciones usar. Construyan
el modelo de forma \textbf{algebraica}, con \texttt{addVars} y \texttt{quicksum}
sobre los conjuntos.
\begin{lstlisting}
import gurobipy as gp
from gurobipy import GRB, quicksum

m = gp.Model("agroandes")

# Ejemplos de sintaxis (nombres genericos, NO son el modelo):
w = m.addVars(Conj1, Conj2, vtype=GRB.BINARY, name="w")    # binaria 0/1
f = m.addVars(Conj1, Conj2, Conj3, lb=0, name="f")         # continua >= 0

m.setObjective(quicksum(...), GRB.MINIMIZE)
m.addConstrs((... for t in T), name="mi_restr")            # nombrela!

m.optimize()
if m.Status == GRB.OPTIMAL:
    print("Optimo:", m.ObjVal)
    print(f[i,j,t].X)   # valor de una variable en la solucion
\end{lstlisting}
Para las restricciones que dependen del periodo anterior (transiciones, balance),
recuerden tratar aparte el \textbf{primer periodo} usando el estado/inventario
inicial dado.
\end{ayuda}

\noindent\rule{\linewidth}{0.4pt}
\newpage
% =============================================================================
\section*{Parte C \quad Dualidad, Precios Sombra y Holgura Complementaria \hfill\normalfont\normalsize(30 pts)}
% =============================================================================

La \textbf{dualidad} y la \textbf{holgura complementaria} son propiedades de los
problemas \textbf{lineales} (LP). Su modelo, en cambio, es un problema
\textbf{entero-mixto} (MIP): tiene variables binarias, y para un MIP los precios
sombra no están definidos de la forma clásica. La técnica estándar para obtener
información dual de un MIP es \textbf{fijar las variables enteras en su valor
óptimo} y resolver el \textbf{LP resultante} (con las binarias ya fijas, el modelo
es lineal en las variables continuas). Sobre ese LP sí valen los precios sombra y
la holgura complementaria.

\begin{ayuda}{Ayuda -- Fijar las binarias y obtener el LP (\texttt{model.fixed})}
\vspace{2pt}
Tras resolver el MIP, \texttt{gurobipy} entrega directamente el \textbf{modelo
fijo}: una copia del modelo con todas las variables enteras fijadas en su valor
óptimo y declaradas continuas. Ese modelo \textbf{sí es un LP} y de él se pueden
leer los atributos duales:
\begin{lstlisting}
m.optimize()                 # 1) resolver el MIP
fx = m.fixed()               # 2) LP con binarias fijas al optimo
fx.optimize()                # 3) resolver el LP
print("Es MIP?", fx.IsMIP)   # -> 0 (es un LP: hay duales disponibles)
\end{lstlisting}
Sobre \texttt{fx} podrán consultar \texttt{.Pi}, \texttt{.Slack} y \texttt{.RC}.
\end{ayuda}

% ── C.1 ──────────────────────────────────────────────────────────────────────
\subsection*{C.1 \quad Precios sombra de la demanda \hfill\normalfont\normalsize(12 pts)}

\begin{enumerate}[label=\alph*), leftmargin=*, topsep=2pt]
    \item Obtenga el \textbf{precio sombra} $\pi_{jt}$ de cada restricción de
          demanda (mercado $j$, periodo $t$) sobre el LP con binarias fijas.
    \item Interprete en \textbf{lenguaje de negocios}: ¿qué significa $\pi_{jt}$?
          ¿Cuánto costaría atender una tonelada adicional de pedido en el mercado
          $X$ en el periodo $t=4$?
    \item Ordene los pares (mercado, periodo) según su precio sombra.
          ¿En qué mercado y periodo es \textbf{más caro} servir una tonelada
          adicional? Explique por qué —relacionándolo con la ventana de cosecha,
          el costo de proceso creciente y la perecibilidad.
\end{enumerate}

\begin{ayuda}{Ayuda -- Precios sombra y holguras en Gurobi}
\begin{lstlisting}
for c in fx.getConstrs():
    print(f"{c.ConstrName:<20} Pi = {c.Pi:8.3f}   Slack = {c.Slack:8.3f}")
\end{lstlisting}
\texttt{c.Pi} es el valor dual (precio sombra) de la restricción: el cambio en el
óptimo por unidad de su lado derecho (RHS). \texttt{c.Slack} es su holgura. Para
que los nombres sean legibles, \textbf{nombren} las restricciones al crearlas
(por ejemplo \texttt{name="demanda"}); así \texttt{ConstrName} las identifica.
\emph{Convención de signos (problema de minimización):} una restricción de
demanda ($\ge$) tiene $\pi\ge 0$ (relajar la exige más costo); una de capacidad
($\le$) tiene $\pi\le 0$ (más recurso reduce el costo).
\end{ayuda}

% ── C.2 ──────────────────────────────────────────────────────────────────────
\subsection*{C.2 \quad Holgura complementaria \hfill\normalfont\normalsize(12 pts)}

La \textbf{holgura complementaria} afirma que, en el óptimo, para cada restricción
se cumple $\pi_i \cdot (\text{holgura}_i)=0$: o la restricción está \textbf{activa}
(holgura cero) o su precio sombra es cero (o ambos). Análogamente, para cada
variable, $(\text{costo reducido})\cdot(\text{valor de la variable})=0$.

\begin{enumerate}[label=\alph*), leftmargin=*, topsep=2pt]
    \item \textbf{Verifique numéricamente} la holgura complementaria sobre el LP
          con binarias fijas: para \textbf{todas} las restricciones, compruebe que
          $|\pi_i \cdot \text{Slack}_i|$ es prácticamente cero (a tolerancia
          numérica). Reporte el máximo valor encontrado.
    \item Elija \textbf{dos} restricciones de demanda con holgura cero y
          \textbf{dos} restricciones con holgura positiva (por ejemplo, de
          capacidad de almacenamiento o de proceso). Para cada una, muestre su par
          $(\pi_i,\text{Slack}_i)$ y explique en palabras por qué el producto es
          cero.
    \item Tome una variable de flujo que valga cero en el óptimo y verifique que
          su \textbf{costo reducido} (\texttt{.RC}) es coherente con la holgura
          complementaria. Interprete qué significa económicamente que esa ruta
          \emph{no} se use.
\end{enumerate}

\begin{ayuda}{Ayuda -- Costos reducidos y verificación de la holgura complementaria}
\begin{lstlisting}
# Costo reducido de las variables (sobre el LP fijo)
for v in fx.getVars():
    if abs(v.X) < 1e-9:              # variables en cero
        print(f"{v.VarName:<14} X = {v.X:6.2f}  RC = {v.RC:8.3f}")

# Verificacion global de holgura complementaria
peor = max(abs(c.Pi * c.Slack) for c in fx.getConstrs())
print("max |Pi * Slack| =", peor)    # debe ser ~ 0
\end{lstlisting}
\texttt{v.RC} es el costo reducido de la variable \texttt{v}. En un óptimo, una
variable con valor cero tiene costo reducido $\ge 0$ (en minimización): usarla
\emph{empeoraría} el costo, por eso no se usa.
\end{ayuda}

% ── C.3 ──────────────────────────────────────────────────────────────────────
\subsection*{C.3 \quad El precio de la perecibilidad \hfill\normalfont\normalsize(6 pts)}

\begin{enumerate}[label=\alph*), leftmargin=*, topsep=2pt]
    \item Identifique cuáles restricciones de \textbf{perecibilidad} están
          \textbf{activas} (holgura cero) en el óptimo y cuál es su precio sombra.
    \item Interprete económicamente ese precio sombra: ¿qué le está diciendo a la
          empresa sobre el valor de que la fruta \emph{durara un poco más} en
          cámara? Conecte su respuesta con la Parte~D.
\end{enumerate}

\noindent\rule{\linewidth}{0.4pt}
\newpage
% =============================================================================
\section*{Parte D \quad El Valor de la Vida Útil \hfill\normalfont\normalsize(15 pts)}
% =============================================================================

La perecibilidad ($L$) no es un dato inamovible: invertir en mejor tecnología de
frío, atmósfera controlada o mejores variedades puede \textbf{extender} la vida
útil de la fruta. Esta parte cuantifica cuánto vale eso.

\begin{enumerate}[label=\alph*), leftmargin=*, topsep=2pt]
    \item Resuelva el modelo (el MIP completo) para \textbf{cada valor} de la vida
          útil $L\in\{1,2,3,4\}$. Para cada $L$ registre si el problema es
          \textbf{factible} y, si lo es, el \textbf{costo total óptimo}.
    \item Grafique el \textbf{costo óptimo} en función de $L$. Comente la forma de
          la curva: ¿qué ocurre con $L=1$? ¿a partir de qué valor de $L$ el costo
          deja de mejorar (se \textbf{aplana})? ¿Por qué?
    \item Calcule el \textbf{ahorro} que produce pasar de $L=2$ a $L=3$. Relacione
          ese ahorro con el precio sombra de la perecibilidad obtenido en la
          Parte~C.3: ¿son consistentes?
    \item En términos de negocio: si extender la vida útil de $2$ a $3$ periodos
          costara una cierta inversión por temporada, ¿cómo usaría este análisis
          para decidir si conviene? ¿Y pasar de $3$ a $4$?
\end{enumerate}

\begin{ayuda}{Ayuda -- Resolver el modelo variando un parámetro}
\vspace{2pt}
Conviene encapsular la construcción y resolución del modelo en una \textbf{función}
que reciba $L$ y devuelva el estado y el costo, para llamarla en un bucle:
\begin{lstlisting}
def resolver(L):
    m = construir_modelo(L)     # arma el MIP con esa vida util
    m.Params.OutputFlag = 0     # silenciar salida dentro del bucle
    m.optimize()
    if m.Status == GRB.OPTIMAL:
        return m.ObjVal
    return None                 # infactible u otro estado

for L in [1, 2, 3, 4]:
    val = resolver(L)
    print(L, "infactible" if val is None else round(val, 1))
\end{lstlisting}
Si para algún $L$ el modelo resulta \textbf{infactible}, \texttt{m.Status} no será
\texttt{GRB.OPTIMAL} (será \texttt{GRB.INFEASIBLE}); repórtenlo como tal en lugar
de leer \texttt{m.ObjVal}.
\end{ayuda}

\noindent\rule{\linewidth}{0.4pt}

% =============================================================================
\section*{Entregables}
% =============================================================================

\subsection*{1.\quad Informe (PDF)}

Estructura mínima requerida (cada sección comienza en una nueva página):

\begin{enumerate}[leftmargin=*, topsep=2pt]
    \item \textbf{Portada:} nombre del curso, Tarea~4, integrantes, profesor,
          ayudante y fecha.
    \item \textbf{Introducción:} descripción del problema de negocio y su
          relevancia.
    \item \textbf{Marco teórico:} conceptos clave (variables binarias de
          estado/transición, big-M, balance de inventario, \textbf{dualidad},
          \textbf{precios sombra} y \textbf{holgura complementaria}).
    \item \textbf{Desarrollo (metodología):} la formulación completa de la
          Parte~A y cómo se implementó y resolvió en \texttt{gurobipy}, incluyendo
          la técnica de fijar binarias para el análisis dual.
    \item \textbf{Resultados:} presente \textbf{solo los resultados importantes}
          —solución óptima, plan de apertura/cierre, inventario, precios sombra,
          verificación de holgura complementaria y la curva de la Parte~D.
          \begin{quote}
          \textit{Ojo:} en esta sección van los resultados puros, sin
          interpretación.
          \end{quote}
    \item \textbf{Análisis de resultados:} interpretación en lenguaje de negocios
          de la solución, de los precios sombra y del valor de la vida útil.
    \item \textbf{Conclusiones:} sea \textbf{específico}. Evite frases genéricas
          del tipo ``el modelo funciona''.
    \item \textbf{Referencias} (si aplica).
\end{enumerate}

\begin{ayuda}{Formalidades del informe (se descuentan puntos por incumplimiento)}
\begin{itemize}[topsep=2pt]
    \item Lenguaje técnico y formal, sin errores ortográficos ni gramaticales.
    \item Texto \textbf{justificado}, fuente tamaño \textbf{12} (Times New Roman,
          Arial u otra tipografía estándar).
    \item Cada figura, gráfico o tabla debe llevar \textbf{título y fuente}. Si el
          material es de elaboración propia, debe indicarse explícitamente.
    \item Sean \textbf{breves y directos}. Se prefiere un análisis corto, directo
          y correcto por sobre páginas de relleno.
\end{itemize}
\end{ayuda}

\subsection*{2.\quad Cuaderno Jupyter (.ipynb)}

\begin{itemize}[topsep=2pt]
    \item Implementación \textbf{comentada} del modelo en \texttt{gurobipy}
          (Parte~B), con variables binarias y continuas.
    \item Extracción e impresión de la solución: plan de apertura/cierre,
          inventario, flujos y verificación de factibilidad.
    \item Análisis dual (Parte~C): fijar binarias con \texttt{model.fixed()},
          precios sombra, costos reducidos y verificación de holgura complementaria.
    \item Análisis del valor de la vida útil (Parte~D) con su gráfico.
    \item Código limpio, ordenado y bien estructurado.
\end{itemize}

\medskip
\noindent\textbf{Formato de entrega:} comprima el informe y el notebook en un
archivo \texttt{.zip} o \texttt{.rar} con el nombre \texttt{Tarea4\_GrupoN}.
\textbf{No entregue los archivos por separado.}

\noindent\rule{\linewidth}{0.4pt}

% =============================================================================
\section*{Plazos y Evaluación}
% =============================================================================

\begin{itemize}[topsep=2pt]
    \item \textbf{Plazo de entrega:} Miercoles 12 de Agosto (conversable).
    \item \textbf{Grupos:} 2 a 3 integrantes (ni más ni menos).
    \item \textbf{Inscripción:} anótense en los grupos de Canvas
          (Personas~$\to$ Grupos~$\to$ Tarea~4\,--\,Grupo~n).
\end{itemize}

\vspace{0.5em}

\begin{center}
\renewcommand{\arraystretch}{1.2}
\begin{tabular}{lc}
\toprule
\textbf{Criterio} & \textbf{Ponderación} \\
\midrule
Parte A: Formulación del modelo                          & 30\,\% \\
Parte B: Implementación y resolución en Gurobi           & 25\,\% \\
Parte C: Dualidad, precios sombra y holgura complementaria & 30\,\% \\
Parte D: El valor de la vida útil                        & 15\,\% \\
\midrule
\textit{Calidad del informe y profundidad del análisis} &
\textit{transversal} \\
\bottomrule
\end{tabular}
\captionof{table}{Criterios de evaluación.}
\end{center}

\vspace{0.3em}
\noindent\textit{La calidad del informe, la profundidad del análisis y la correcta
interpretación de los resultados en lenguaje de negocios impactan de forma
transversal en todas las partes.}

\noindent\rule{\linewidth}{0.4pt}

% =============================================================================
\section*{Notas Importantes}
% =============================================================================

\begin{itemize}[topsep=2pt]
    \item \textbf{Deben usar \texttt{gurobipy}} para la implementación.
    \item Cada variable debe declararse con su \textbf{dominio y tipo correctos}
          (continua o binaria, según lo que represente). Identificar los tipos es
          parte de la evaluación.
    \item \textbf{Nombren} las restricciones al crearlas: es indispensable para
          leer los precios sombra y las holguras en la Parte~C.
    \item Construyan el modelo de forma \textbf{algebraica} (con conjuntos,
          \texttt{addVars} y \texttt{quicksum}); no escriban los números uno a uno
          dentro de las restricciones.
    \item Los atributos duales (\texttt{.Pi}, \texttt{.Slack}, \texttt{.RC}) solo
          están disponibles \textbf{sobre el LP con binarias fijas}
          (\texttt{model.fixed()}) y solo si el estado es \texttt{GRB.OPTIMAL}.
    \item En el informe y en el código deben figurar los \textbf{nombres de todos
          los integrantes} del grupo.
\end{itemize}

\vspace{1em}
\begin{center}
\textit{Consultas al correo
\href{mailto:vramireza@udd.cl}{\texttt{vramireza@udd.cl}} o al
\texttt{+56\,9\,9912\,1814}
(respondo siempre que puedo, dentro de horarios razonables).}
\end{center}

% =============================================================================
\end{document}
% =============================================================================
